\documentclass[aspectratio=169,11pt]{beamer}
\usepackage{fontspec}
\usepackage{listings}
\usepackage{xcolor}
\usepackage{booktabs}
\usepackage{array}

% ================= Palette (identique au livre Tome 1) =================
\definecolor{PyBleu}{HTML}{1A4A7A}
\definecolor{PyBleuMoyen}{HTML}{2D7AB3}
\definecolor{PyBleuClair}{HTML}{EAF1F8}
\definecolor{PyJaune}{HTML}{FFD43B}
\definecolor{PyNuit}{HTML}{142840}
\definecolor{PyTexte}{HTML}{1F2A37}
\definecolor{PyGris}{HTML}{5A6B7E}
\definecolor{PyRouge}{HTML}{B3442D}
\definecolor{PyVert}{HTML}{2E7D5B}
\definecolor{CodeTexte}{HTML}{E8EEF6}

% ================= Code Python =================
\lstdefinestyle{py}{
  language=Python,
  keywordstyle=\color{PyJaune}\bfseries,
  stringstyle=\color{PyJaune},
  commentstyle=\color{PyGris}\itshape,
  showstringspaces=false,
  breaklines=true,
  backgroundcolor=\color{PyNuit},
  basicstyle=\ttfamily\footnotesize\color{CodeTexte},
  columns=fullflexible,
  keepspaces=true,
  xleftmargin=6pt,
  framexleftmargin=6pt,
}
\lstdefinestyle{pysmall}{
  language=Python,
  keywordstyle=\color{PyJaune}\bfseries,
  stringstyle=\color{PyJaune},
  commentstyle=\color{PyGris}\itshape,
  showstringspaces=false,
  breaklines=true,
  backgroundcolor=\color{PyNuit},
  basicstyle=\ttfamily\scriptsize\color{CodeTexte},
  columns=fullflexible,
  keepspaces=true,
  xleftmargin=5pt,
  framexleftmargin=5pt,
}

% ================= Titres de diapositives =================
\setbeamercolor{structure}{fg=PyBleu}
\setbeamerfont{frametitle}{series=\bfseries, size=\Large}
\setbeamercolor{frametitle}{bg=PyBleu, fg=white}
\setbeamertemplate{frametitle}{%
  \nointerlineskip
  \begin{beamercolorbox}[wd=\textwidth, ht=2.1ex, dp=0.9ex,
                          leftskip=0.35cm, rightskip=0.35cm]{frametitle}%
    \usebeamerfont{frametitle}\insertframetitle%
  \end{beamercolorbox}%
  \vspace*{-0.4pt}%
  {\color{PyJaune}\rule{\textwidth}{1.6pt}}%
}

% ================= Pied de page =================
\setbeamerfont{footline}{size=\tiny}
\setbeamercolor{footline}{bg=white, fg=PyGris}
\newcommand{\PiedDePage}{Python en 24 heures \textbar\ Séance 07 : Closures, chaînes et fichiers}
\setbeamertemplate{footline}{%
  \leavevmode%
  \hbox{%
    \begin{beamercolorbox}[wd=\paperwidth, ht=2.2ex, dp=1.1ex,
                            leftskip=0.35cm, rightskip=0.35cm]{footline}%
      \usebeamerfont{footline}\PiedDePage
      \hfill\insertframenumber\,/\,\inserttotalframenumber%
    \end{beamercolorbox}}%
}

\setbeamertemplate{navigation symbols}{}
\setbeamertemplate{itemize item}{\color{PyBleuMoyen}$\blacktriangleright$}
\setbeamertemplate{itemize subitem}{\color{PyBleuMoyen}\textbullet}

% ================= Blocs pedagogiques =================
\setbeamercolor{block title}{bg=PyBleu, fg=white}
\setbeamercolor{block body}{bg=PyBleuClair, fg=PyTexte}
\setbeamercolor{block title alerted}{bg=PyRouge, fg=white}
\setbeamercolor{block body alerted}{bg=PyRouge!9, fg=PyTexte}
\setbeamercolor{block title example}{bg=PyJaune, fg=PyNuit}
\setbeamercolor{block body example}{bg=PyJaune!14, fg=PyTexte}

% ================= Metadonnees =================
\hypersetup{
  pdftitle={Python en 24 heures : Séance 07 : Closures, chaînes et fichiers},
  pdfauthor={HOUNKPE Mawulikplimi Roland et DJABON Ounimborbitibou},
  pdfsubject={De l'algorithmique à la programmation orientée objet},
  pdfkeywords={Python, algorithmique, POO, formation, exercices corrigés}
}

\begin{document}

{%
\setbeamercolor{background canvas}{bg=PyNuit}
\begin{frame}[plain]
\vspace{10mm}
{\color{PyJaune}\small\bfseries FORMATION « PYTHON EN 24 HEURES » \quad\textbar\quad SÉANCE 07\par}
\vspace{4mm}
{\color{white}\Huge\bfseries Fonctions avancées : closures, décorateurs, chaînes et fichiers\par}
\vspace{4mm}
{\color{PyBleuClair}\large Chapitres 13 et 14 du Tome 1\par}
\vspace{7mm}
{\color{PyJaune}\rule{34mm}{2.2pt}\par}
\vspace{7mm}
{\color{white}\normalsize 1 h de cours (présentation Beamer) \quad$\bullet$\quad 1 h de pratique (notebook Jupyter)\par}
\vspace{2.5mm}
{\color{PyBleuClair}\normalsize HOUNKPE Mawulikplimi Roland et DJABON Ounimborbitibou\par}
\end{frame}
}
\begin{frame}{Objectifs de la séance}
\begin{itemize}
  \item Comprendre la portée LEGB et l'imbrication de fonctions
  \item Écrire des closures : des fonctions qui gardent en mémoire leur contexte
  \item Comprendre et écrire un décorateur avec @
  \item Maîtriser les méthodes essentielles des chaînes : strip, lower, split, join
  \item Lire et écrire des fichiers proprement avec with open
\end{itemize}
\end{frame}
\begin{frame}{Rappel de la séance 06}
Une fonction est un objet comme un autre : on peut la retourner, la passer en paramètre et même la fabriquer à la chaîne. C'est la clé des closures.

\begin{itemize}
  \item Dict : accès par clé, .items(), schéma de comptage
  \item Set : unicité, appartenance rapide, union et intersection
  \item Fonctions : def, paramètres, valeurs par défaut, return
  \item Portée : une variable créée dans une fonction disparaît à la sortie
  \item Choisir la structure : position ? list. Nom ? dict. Unicité ? set
\end{itemize}
\end{frame}
\begin{frame}{La règle LEGB : où Python cherche un nom}
\renewcommand{\arraystretch}{1.3}\small
\begin{center}
\begin{tabular}{lll}
\toprule
\textbf{\color{PyBleu}Niveau} & \textbf{\color{PyBleu}Signification} & \textbf{\color{PyBleu}Durée de vie} \\
\midrule
L & Local : dans la fonction courante & à chaque appel \\
E & Englobante : la fonction qui contient & tant que la closure vit \\
G & Globale : au niveau du module & toute l'exécution \\
B & Built in : print, len, sum... & toujours \\
\bottomrule
\end{tabular}
\end{center}

\vspace{4pt}

{\small\color{PyGris}Python cherche de bas en haut et s'arrête au premier nom trouvé. Une affectation crée un nom LOCAL, sauf déclaration nonlocal ou global.}
\end{frame}

\begin{frame}[fragile]{Fonctions imbriquées}

\begin{lstlisting}[style=py]
def exterior(x):
    def interieur(y):
        return x + y      # x vient de la portée englobante
    return interieur

ajoute_10 = exterior(10)
print(ajoute_10(5))       # 15
print(ajoute_10(32))      # 42
\end{lstlisting}
\vspace{0.5pt}

{\small\color{PyGris}\itshape La fonction interne peut lire les variables de la fonction externe. Retourner `interieur` exporte la fonction, avec son souvenir de x.}

\end{frame}

\begin{frame}{La closure : la fonction à mémoire}
\begin{itemize}
  \item Une \textbf{closure} est une fonction interne retournée qui \textbf{conserve} les variables de sa portée de naissance
  \item Chaque appel de la fabrique crée une closure \textbf{indépendante}
  \item `ajoute\_10.\_\_closure\_\_` révèle la valeur retenue
  \item Usage : fabriques de fonctions configurées, compteurs, caches
  \item C'est le socle des décorateurs, vus dans trois diapositives
\end{itemize}
\end{frame}

\begin{frame}[fragile]{Fabrique de multiplicateurs}

\begin{lstlisting}[style=py]
def fabrique_multiplicateur(factor):
    def multiplier(valeur):
        return valeur * factor   # factor est retenu
    return multiplier

double = fabrique_multiplicateur(2)
triple = fabrique_multiplicateur(3)
print(double(7), triple(7))   # 14 21
\end{lstlisting}
\vspace{0.5pt}

{\small\color{PyGris}\itshape double et triple sont deux closures distinctes : chacune a retenu son propre factor. Changer double ne touche pas triple.}

\end{frame}

\begin{frame}[fragile]{Modifier la variable englobante : nonlocal}
\begin{columns}[t]
\begin{column}{0.42\textwidth}
\begin{itemize}
  \setlength{\itemsep}{1.5pt}\setlength{\parskip}{0pt}
  \item Lire la variable englobante : automatique
  \item La \textbf{réaffecter} : il faut déclarer `nonlocal compteur`
  \item Sans nonlocal, l'affectation créerait une variable locale de même nom
  \item Schéma classique : le compteur vit dans la closure, personne d'autre ne peut y toucher
\end{itemize}
\end{column}
\hspace{0.02\textwidth}
\begin{column}{0.55\textwidth}
\begin{lstlisting}[style=pysmall]
def fabrique_compteur():
    compteur = 0
    def incrementer():
        nonlocal compteur
        compteur += 1
        return compteur
    return incrementer

prochain = fabrique_compteur()
print(prochain(), prochain(), prochain())
print(prochain.__closure__[0].cell_contents)
\end{lstlisting}
\end{column}
\end{columns}
\end{frame}

\begin{frame}[fragile]{Le décorateur : une closure appliquée}

\begin{lstlisting}[style=py]
def majuscules(fonction):
    def emballage(*args, **kwargs):
        resultat = fonction(*args, **kwargs)
        return resultat.upper()
    return emballage

@majuscules
def salutation(nom):
    return f"bonjour {nom}"

print(salutation("Awa"))   # BONJOUR AWA
\end{lstlisting}
\vspace{0.5pt}

{\small\color{PyGris}\itshape @majuscules est un raccourci pour : salutation = majuscules(salutation). La fonction d'origine vit toujours, emballée.}

\end{frame}

\begin{frame}[fragile]{Préserver l'identité : functools.wraps}
\begin{columns}[t]
\begin{column}{0.42\textwidth}
\begin{itemize}
  \setlength{\itemsep}{1.5pt}\setlength{\parskip}{0pt}
  \item Après emballage, `salutation.\_\_name\_\_` vaudrait « emballage »
  \item @wraps(fonction) recopie nom, docstring et métadonnées
  \item Indispensable pour le débogage et la documentation
  \item Les arguments *args et **kwargs rendent le décorateur universel
\end{itemize}
\end{column}
\hspace{0.02\textwidth}
\begin{column}{0.55\textwidth}
\begin{lstlisting}[style=pysmall]
from functools import wraps

def compte_appels(fonction):
    @wraps(fonction)
    def emballage(*args, **kwargs):
        emballage.appels += 1
        return fonction(*args, **kwargs)
    emballage.appels = 0
    return emballage
\end{lstlisting}
\end{column}
\end{columns}
\end{frame}

\begin{frame}[fragile]{Les méthodes des chaînes}

\begin{lstlisting}[style=py]
phrase = "  Bonjour la Classe  "
print(phrase.strip())            # retire les espaces aux bords
print(phrase.strip().lower())    # tout en minuscules
mots = "riz,yam,huile".split(",")
print(mots)                      # ['riz', 'yam', 'huile']
print(" + ".join(mots))          # riz + yam + huile
print("abc".replace("a", "X"))   # Xbc
\end{lstlisting}
\vspace{0.5pt}

{\small\color{PyGris}\itshape Les chaînes sont immuables : chaque méthode retourne une NOUVELLE chaîne. Enchaînez les comme des filtres : phrase.strip().lower().split().}

\end{frame}

\begin{frame}[fragile]{Fichiers : la règle with open}
\begin{columns}[t]
\begin{column}{0.42\textwidth}
\begin{itemize}
  \setlength{\itemsep}{1.5pt}\setlength{\parskip}{0pt}
  \item `open(nom, mode)` ouvre : « r » lire, « w » écrire en écrasant, « a » ajouter à la fin
  \item \textbf{with} ferme le fichier même en cas d'erreur : jamais de fichier oublié
  \item `.read()` lit tout, `.readlines()` une liste de lignes, boucler lit ligne à ligne
  \item Chaque ligne lue se termine par \textbackslash\{\}n : `strip()` la nettoie
  \item On détaillera les erreurs de fichier à la séance 11
\end{itemize}
\end{column}
\hspace{0.02\textwidth}
\begin{column}{0.55\textwidth}
\begin{lstlisting}[style=pysmall]
with open("notes.txt", "w") as f:
    f.write("12\n15\n9\n")

with open("notes.txt", "r") as f:
    total = 0
    for ligne in f:
        total += int(ligne.strip())
print(total)   # 36
\end{lstlisting}
\end{column}
\end{columns}
\end{frame}

\begin{frame}{Astuce}
\begin{exampleblock}{Explorateur de closure}
Après avoir créé une closure, inspectez la : `ma\_fonction.\_\_closure\_\_` est un tuple de cellules, et `.cell\_contents` montre la valeur retenue. Ce microscope rend concret un concept souvent abstrait : la fonction garde bien une boîte avec la variable dedans.
\end{exampleblock}
\end{frame}

\begin{frame}{Attention}
\begin{alertblock}{w écrase, a ajoute}
Ouvrir un fichier en mode « w » efface immédiatement tout son contenu, avant même la première écriture. Pour enrichir un journal au fil des jours, utilisez « a ». Et rappelez vous que les chaînes sont immuables : `phrase.upper()` ne modifie rien si vous n'affectez pas le résultat.
\end{alertblock}
\end{frame}

\begin{frame}{À retenir}
\begin{block}{L'essentiel de la séance 07}
LEGB régit la recherche des noms ; la closure retient la portée englobante et nonlocal autorise sa modification ; le décorateur enveloppe une fonction pour lui ajouter un comportement. Les chaînes se traitent par chaînage de méthodes immuables. Les fichiers se manipulent toujours avec with open : lire, écrire, ajouter.
\end{block}
\end{frame}

\begin{frame}{Travaux pratiques (1 h) : Closures, décorateurs, chaînes et fichiers}
Cinq exercices qui font dialoguer fonctions avancées et données textuelles. Le dernier exercice écrit réellement un fichier sur le disque, puis le supprime proprement.

\begin{itemize}
  \setlength{\itemsep}{1.5pt}\setlength{\parskip}{0pt}
  \item \textbf{Exercice 1 (Basique)} : Nettoyage de chaîne
  \item \textbf{Exercice 2 (Basique)} : Palindrome
  \item \textbf{Exercice 3 (Intermédiaire)} : Fabrique de puissances (closure)
  \item \textbf{Exercice 4 (Intermédiaire)} : Décorateur compteur d'appels
  \item \textbf{Exercice 5 (Défi)} : Journal de bord dans un fichier
\end{itemize}
\vspace{4pt}

{\small\color{PyGris}Ouvrez le notebook « Seance\_07\_closures\_decorateurs\_fichiers.ipynb » pour commencer.}
\end{frame}
\begin{frame}{Pour la prochaine séance}
\begin{exampleblock}{Conseil de travail}
Séance 08 : la programmation orientée objet. Tout ce que vous avez écrit jusqu'ici était procédural ; nous allons ranger données et comportements dans des classes. Relisez le schéma de comptage : la classe Compte en aura besoin pour son historique.
\end{exampleblock}
\end{frame}

\end{document}
